\scanpage{175}\mspage{120}
du diagramme 76, à la seule exception de $\mathrm{Cat}_{/S}$ elle-même. Si $X$ est dans $\mathrm{Cat}_{/S}$, sans être ni lisse ni Hot-fibrée sur $S$, comment trouver \unsure{l'|son }image inverse dans $\mathrm{HOT}_W(S')$ ? On a le choix entre $\unsure{\Phi|Q}$ et $C$ pour améliorer la situation ! On sait que $X$ définit la $=$ chose dans $\mathrm{HOT}_W(S)$ que $\Phi(X),C(X)$, qui sont l'une lisse (et $=$ fibrée scindée), l'autre Hot-fibrée\deleted{\unsure{s}} (et même parfaite) sur $S$. Donc

\uline{Corollaire 2} Si $X$ dans $\mathrm{Cat}_{/S}$ est quelconque, alors
\[
f^*(\mathrm{hot}_{W,S}(X))\simeq\mathrm{hot}_{W,S'}(\Phi(X)\times_S S')\simeq\mathrm{hot}_{W,S'}(C(X)\times_S S').
\]

On peut s'attendre, d'autre part, que pour certains changements de base particuliers, la formule de la prop.~1 soit valable sans restriction sur $X$ dans $\mathrm{Cat}_{/S}$. Il doit revenir $\pm$ au même que le foncteur image inverse
\[
X\longmapsto X\times_S S'
\]
est compatible avec les localiseurs $W_S^d,W_{S'}^d$.

La \deleted{\unsure{con}} condition de $W$-équivalence ``colocale sur $S$'' doit être vue comme une condition de nature ``homologique'', pour le dérivateur $\mathrm{HOT}$. Elle s'associe, pour un dérivateur quelconque $\mathbb D$, à la condition (pour $f:X\longrightarrow Y$ de $\mathrm{Cat}_{/S}$) que \deleted{les} les $ {}_{s\backslash}X\unsure{\text{---}|\longrightarrow}{}_{s\backslash}Y$ soient des $\mathbb D$-équivalences,
% archive sheet number: 174
